\documentclass[11pt]{article}
\usepackage[margin=1in]{geometry}
\usepackage{amsmath,amssymb,amsthm}
\usepackage{graphicx}
\usepackage[hidelinks]{hyperref}

\newtheorem{theorem}{Theorem}
\newtheorem{lemma}{Lemma}
\newcommand{\Z}{\mathbb Z}
\newcommand{\R}{\mathbb R}
\newcommand{\C}{\mathbb C}
\newcommand{\Hcal}{\mathcal H}
\newcommand{\Ss}{\mathsf S}
\newcommand{\E}{\mathsf E}
\newcommand{\one}{\mathbf 1}
\DeclareMathOperator{\spanop}{span}

\title{Odd Discrete Tori in Trigonometric Chaos}
\author{Run \texttt{fa\_banach\_001}, agent \texttt{agent\_lane\_04}}
\date{2026-07-18}

\begin{document}
\maketitle

\section*{Status}

This packet gives a likely valid full answer to the extracted signal in
Remark 2.4 of Cano-Marmol, Conde-Alonso, and Parcet
\cite{CCP2022}: an explicit orthonormal basis for the cocycle Hilbert
space of the odd cyclic product $\Z_{2m+1}^n$, and the resulting
balanced Fourier-truncation inequality.  The result is a completion of a
source-paper computation left to the interested reader, not a claim that a
separate major conjecture is resolved.

\section*{Source Signal}

The source paper proves the even-cardinal discrete-torus case in Section 2.2
and then records the following odd-cardinal variant in Remark 2.4 on page 16.

\begin{center}
\includegraphics[width=\textwidth]{figures/open_problem_crop.png}
\end{center}

In words, the authors note that for $\Z_{2m+1}^n$ the word length
\[
  \psi(g)=\sum_{j=1}^n |g_j|,\qquad |a|=\min\{a,2m+1-a\},
\]
is conditionally negative, that the associated cocycle Hilbert space has
dimension $2mn$, and that Theorems A and B apply; what is left open in the
remark is an explicit orthonormal basis.

\section*{Proof Intuition}

For the even cycle, the source basis is made of oriented edges, but opposite
edges cancel in pairs.  In an odd cycle there is no antipodal pairing, so the
edge basis is no longer orthogonal in the same elementary way.  The clean
replacement is to diagonalize the cycle distance by the nonzero Fourier modes
of $\Z_{2m+1}$.  The negative Fourier coefficients of the distance give the
normalizing constants.  Taking real and imaginary parts produces $2m$ real
orthonormal vectors for each coordinate.  These coordinate Fourier planes are
rotated by the cocycle action, so they give the same admissible coordinate
splitting used by the source paper.  The coordinate support derivative
$\partial_j\lambda(g)=\one_{\{g_j\ne0\}}\lambda(g)$ then absorbs every
cocycle derivative in the $j$th coordinate, which is exactly the distinguished
derivative hypothesis of Theorem B in \cite{CCP2022}.

\section*{Theorem}

Let $N=2m+1$ and write elements of $\Z_N$ as $0,1,\ldots,N-1$, with arithmetic
modulo $N$.  Put
\[
  |a|=\min\{a,N-a\},\qquad
  \psi(g)=\sum_{j=1}^n |g_j|,\qquad g\in \Z_N^n.
\]
Let $\Hcal_\psi$ be the real cocycle Hilbert space obtained from the Gromov
form
\[
  \langle \delta_g,\delta_h\rangle_\psi
  =\frac{1}{2}\bigl(\psi(g)+\psi(h)-\psi(g-h)\bigr).
\]
For $1\le r\le m$ set $\theta_r=2\pi r/N$ and
\[
  q_r
  =-\sum_{a=0}^{N-1}|a|\cos(\theta_r a)
  =\frac{1-(-1)^r\cos(\pi r/N)}{2\sin^2(\pi r/N)}  >0 .
\]
For $1\le j\le n$, define
\[
  c_{j,r}=-\frac{2}{\sqrt{Nq_r}}\sum_{a=0}^{N-1}
       \cos(\theta_r a)\,\delta_{ae_j},
  \qquad
  s_{j,r}= \frac{2}{\sqrt{Nq_r}}\sum_{a=0}^{N-1}
       \sin(\theta_r a)\,\delta_{ae_j}.
\]
Then
\[
  \bigl\{c_{j,r},s_{j,r}: 1\le j\le n,\ 1\le r\le m\bigr\}
\]
is an orthonormal basis of $\Hcal_\psi$.

Moreover, if $\Hcal_{\psi,j}=\spanop\{c_{j,r},s_{j,r}:1\le r\le m\}$, then
$\Hcal_\psi=\bigoplus_{j=1}^n\Hcal_{\psi,j}$ is invariant under the cocycle
action.  The conditional expectations over the coordinate subgroups
$\Z_N^\Ss$ form an admissible family of Fourier truncations.  If
\[
  \partial_j\lambda(g)=\one_{\{g_j\ne0\}}\lambda(g),
\]
then $\{\partial_j\}_{j=1}^n$ is a distinguished family of derivatives.  Hence,
for every $p\ge2$, every $1\le k\le n$, and every mean-zero
$f:\Z_N^n\to\C$,
\[
 \frac{1}{\binom{n}{k}}\sum_{\substack{\Ss\subseteq[n]\\|\Ss|=k}}
 \left\|\sum_{g\in\Z_N^\Ss}\widehat f(g)
       \exp\left(\frac{2\pi i}{N}\langle\cdot,g\rangle\right)\right\|_p^p
 \lesssim_p
 \frac{k}{n}\sum_{j=1}^n \|\partial_j f\|_p^p
 +\left(\frac{k}{n}\right)^{p/2}\|f\|_p^p .
\]

\section*{Proof}

We first record the zero-sum form of the Gromov inner product.  If
$x=\sum_g a_g\delta_g$ and $y=\sum_h b_h\delta_h$ satisfy
$\sum_g a_g=\sum_h b_h=0$, then
\begin{equation}\label{eq:zero-sum}
  \langle x,y\rangle_\psi
  =-\frac{1}{2}\sum_{g,h}a_gb_h\,\psi(g-h).
\end{equation}
Indeed, the two terms containing $\psi(g)$ and $\psi(h)$ vanish after summing
against the zero total coefficients.

Let $\kappa(a)=|a|$ on $\Z_N$.  For $1\le r\le m$ its nonzero Fourier
coefficient is
\begin{equation}\label{eq:q-fourier}
  \widehat\kappa(r)=\sum_{a=0}^{N-1}\kappa(a)e^{-i\theta_r a}
  =2\sum_{a=1}^m a\cos(\theta_r a)=-q_r .
\end{equation}
The displayed closed form for $q_r$ follows from
\[
  \sum_{a=1}^m az^a=\frac{z\bigl(1-(m+1)z^m+mz^{m+1}\bigr)}{(1-z)^2},
  \qquad z^N=1,\ z\ne1,
\]
because $z^m+z^{m+1}=2(-1)^r\cos(\pi r/N)$ and
$z/(1-z)^2=-1/(4\sin^2(\pi r/N))$.

Now fix one coordinate $j$.  The real functions
$a\mapsto \cos(\theta_r a)$ and $a\mapsto\sin(\theta_r a)$, $1\le r\le m$,
are the real nonzero Fourier modes of $\Z_N$.  By \eqref{eq:zero-sum} and
\eqref{eq:q-fourier},
\[
 \left\|\sum_a \cos(\theta_r a)\delta_{ae_j}\right\|_{\Hcal_\psi}^2
 =
 \left\|\sum_a \sin(\theta_r a)\delta_{ae_j}\right\|_{\Hcal_\psi}^2
 =\frac{Nq_r}{4},
\]
and different real modes are orthogonal.  If $j\ne j'$, the coordinate
supports are also orthogonal: by additivity of $\psi$, every mixed inner
product is a sum of terms containing either the total coefficient of the
$j$-mode or the total coefficient of the $j'$-mode, both of which are zero.
The displayed normalization therefore gives an orthonormal family of
$2mn$ vectors.  Since the source remark identifies
$\dim\Hcal_\psi=2mn$, the family is an orthonormal basis.

The useful coefficient formulas are, for $t=g_j\in\{0,\ldots,N-1\}$,
\begin{equation}\label{eq:coefficients}
  \langle \delta_g,c_{j,r}\rangle_\psi
  =\sqrt{\frac{q_r}{N}}\bigl(1-\cos(\theta_r t)\bigr),\qquad
  \langle \delta_g,s_{j,r}\rangle_\psi
  =\sqrt{\frac{q_r}{N}}\sin(\theta_r t).
\end{equation}
For instance, if $C_r(a)=\cos(\theta_r a)$, then
\[
  \left\langle \delta_g,\sum_a C_r(a)\delta_{ae_j}\right\rangle_\psi
  =\frac{1}{2}\left(\widehat\kappa(r)-\widehat\kappa(r)\cos(\theta_r t)\right)
  =-\frac{q_r}{2}\bigl(1-\cos(\theta_r t)\bigr),
\]
and the sign in the definition of $c_{j,r}$ gives the first formula.  The sine
formula is the same convolution calculation.

The formulas \eqref{eq:coefficients} show in particular that every derivative
in the $j$th coordinate vanishes on Fourier coefficients with $g_j=0$.
Consequently
\[
  \partial_{c_{j,r}}\circ\partial_j=\partial_{c_{j,r}},
  \qquad
  \partial_{s_{j,r}}\circ\partial_j=\partial_{s_{j,r}},
\]
where $\partial_u\lambda(g)=2\pi i\langle\delta_g,u\rangle_\psi\lambda(g)$.
Thus the coordinate support derivatives form a distinguished family once the
coordinate splitting is known to be invariant.

It remains only to note that invariance.  For a zero-sum vector
$x=\sum_a x_a\delta_{ae_j}$, the source cocycle action gives
$\alpha_g x=\sum_a x_a\delta_{g+ae_j}$ in the quotient.  This is translation
by $g_j$ on the $j$th cyclic coordinate.  Hence
\[
  \alpha_g c_{j,r}
  =\cos(\theta_r g_j)c_{j,r}-\sin(\theta_r g_j)s_{j,r},
  \qquad
  \alpha_g s_{j,r}
  =\sin(\theta_r g_j)c_{j,r}+\cos(\theta_r g_j)s_{j,r}.
\]
So each $\Hcal_{\psi,j}$ is $\alpha$-invariant, and the coordinate subgroup
conditional expectations over $\Z_N^\Ss$ are admissible exactly as in the even
case treated in the source paper.

Applying Theorem B of \cite{CCP2022} gives the stated inequality, up to the
extra adjoint derivative terms present in the general theorem.  Since the
group is abelian and the multiplier $\one_{\{g_j\ne0\}}$ is symmetric under
$g\mapsto -g$, we have $\|\partial_j(f^*)\|_p=\|\partial_j f\|_p$; these terms
are absorbed into the implicit constant depending only on $p$.

\section*{Verification Notes}

The script \texttt{code/verify\_odd\_cycle\_onb.py} checks the finite Gram
matrix and coefficient identities for $N=3,5,7,9,11$ and $n=1,2,3$ using the
Gromov form directly.  It is only a numerical sanity check; the proof above is
symbolic.

Cheap-index search found no prior packet for arXiv:2209.05986.  A web search
for the exact odd-cardinal/orthonormal-basis phrases found the source preprint
and its 2024 journal publication, but no separate explicit odd-cycle basis
answer.  The related arXiv:2209.05991 Part II problem has a prior no-promotion
attempt in this run and is mathematically separate.

\section*{Limitations and Human Review Recommendation}

This packet should be reviewed as a compact completion of Remark 2.4, not as a
solution to a broad open conjecture.  The main review points are the sign and
normalization in \eqref{eq:q-fourier} and \eqref{eq:coefficients}, and the
use of Theorem B after verifying the distinguished derivative condition.

\begin{thebibliography}{9}
\bibitem{CCP2022}
Antonio Ismael Cano-Marmol, Jose M. Conde-Alonso, and Javier Parcet.
\newblock Trigonometric chaos and {$X_p$} inequalities I -- Balanced Fourier
truncations over discrete groups.
\newblock arXiv:2209.05986, 2022. Published in Analysis \& PDE 17(7),
2561--2584, 2024.
\end{thebibliography}

\end{document}
